\GalSourcePageBoundary{089}{L0088}{60}{60}
soluble par radicaux, à moins que l'on n'eût $p + 1 = 2^{n}$, $n$~étant
un certain nombre.

Nous pouvons supposer que le groupe ne contienne que des
substitutions de l'ordre~$p$ et de l'ordre~$p + 1$. Toutes les substitutions
de l'ordre~$p + 1$ seront par conséquent semblables, et leur
période sera de deux termes.

Prenons donc l'expression
\[
\biggl(k,\ \frac{mk + n}{rk + s}\biggr),
\]
et voyons dans quel cas cette substitution peut avoir une période
de deux termes. Il faut pour cela que la substitution inverse se
confonde avec elle. La substitution inverse est
\[
\biggl(k,\ \frac{-sk + n}{rk - m}\biggr).
\]

Donc on doit avoir $m = -s$, et toutes les substitutions en
question seront
\[
\biggl(k,\ \frac{mk + n}{k - m}\biggr),
\]
ou encore
\[
k,\ m + \frac{N}{k - m},
\]
$N$~étant un certain nombre qui est le même pour toutes les substitutions,
puisque ces substitutions doivent être transformées les
unes dans les autres par toutes les substitutions de l'ordre~$p$,
$(k, k+m)$; or ces substitutions doivent, de plus, être conjuguées
les unes des autres. Si donc
\[
\biggl(k,\ m + \frac{N}{k - m}\biggr), \qquad
\biggl(k,\ n + \frac{N}{k - n}\biggr)
\]
sont deux pareilles substitutions, il faut que l'on ait
\[
n + \frac{N}{\dfrac{N}{k-m} + m - n} =
m + \frac{N}{\dfrac{N}{k-n} + n - m},
\]
savoir
\[
(m - n)^{2} = 2N.
\]

Donc la différence entre deux valeurs de~$m$ ne peut acquérir
